Observations\section{Background}
\label{sec:background}
\subsection{Generative LLM Inference}

Generative LLMs~\cite{XX} perform autoregressive inference, which typically consists of two serving phases: prefill and decode. In the prefill phase, the model processes the prompt tokens and populates the key-value (KV) cache for all transformer layers. In the decode phase, it generates output tokens iteratively using the cached states from prior tokens. The two phases exhibit distinct execution characteristics. Prefill exposes high parallelism across prompt tokens and is typically compute-bound, whereas decode is inherently sequential and often memory-bound due to repeated KV-cache accesses. Within each phase, execution consists of multiple transformer operators, notably self-attention and FFN operations, which exhibit different compute and memory behaviors.

LLM serving performance is typically evaluated using time to first token (TTFT), time per output token (TPOT), and throughput. TTFT measures the latency from request arrival to the first generated token and is largely determined by prompt processing in prefill, whereas TPOT captures the average per-token latency during the decode phase. TTFT and TPOT capture different aspects of user-perceived latency and are often associated with distinct service-level objectives (SLOs). In production deployments, TTFT and TPOT depend on workload characteristics, system configuration, and runtime state, such as prompt and output lengths, request concurrency, batch size, tensor parallelism, and KV-cache utilization. As a result, consistently meeting these SLOs can be challenging.

\subsection{Prefill/Decode Disaggregation}
\label{sec2.2}

In many LLM serving systems~\cite{XX}, prefill and decode are collocated on the same GPUs and interleaved through continuous batching~\cite{XX}. Although this design can improve hardware utilization and throughput, it also introduces significant inter-phase interference. Prefill processes the full prompt with high parallelism, whereas decode proceeds iteratively and is more sensitive to per-token latency, often bottlenecked by KV-cache accesses. Consequently, interleaving prefill and decode can delay ongoing decode iterations, leading to head-of-line blocking, inflated per-token latency, and difficulty in simultaneously meeting TTFT and TPOT SLOs.

To mitigate this issue, existing work~\cite{XX} has proposed techniques such as chunked prefill, which splits long prompts into smaller chunks to reduce interference with decode. However, such methods only alleviate interference rather than eliminating it, and often require workload-dependent tuning that is difficult to generalize across dynamic workloads. To address this limitation, recent work~\cite{XX} proposes prefill/decode (P/D) disaggregation, which separates prompt processing and autoregressive decoding onto distinct GPU pools or worker groups. This separation reduces cross-phase interference and enables phase-specific batching, provisioning, and scaling, allowing systems to better balance TTFT, TPOT, and throughput under dynamic workloads. However, despite these benefits, P/D disaggregation is a phase-level mechanism and overlooks operator-level differences in compute intensity, memory behavior, and resource demands within each phase, potentially leaving additional optimization opportunities unexplored.

\subsection{DVFS for Energy-Efficient LLM Serving}
DVFS improves energy efficiency by tuning GPU frequency settings to balance performance and energy consumption during inference. In LLM serving, DVFS is particularly relevant because GPU inference is both power-intensive and latency-sensitive. DVFS has been explored in inference workloads to reduce energy consumption while maintaining latency objectives~\cite{XX,XX}. In LLM serving, phase-aware approaches~\cite{XX,XX} use phase-specific DVFS policies for prefill and decode to match their different performance and energy behaviors better, improving uniform DVFS policies by aligning frequency control with phase-level heterogeneity.

However, these works still apply DVFS settings at coarse granularity. As discussed in ~\ref{sec2.2}, each phase consists of multiple operators, such as self-attention and FFNs~\cite{XX}, that differ in compute intensity, memory behavior, and latency sensitivity.As a result, a DVFS setting that is suitable for one operator may be suboptimal for another, even within the same phase. Moreover, the optimal setting for a given operator is not fixed and can vary across runtime conditions, such as prompt length, output length, and request concurrency. Therefore, improving the energy efficiency of LLM serving requires finer-grained and context-aware DVFS control beyond phase-level policies.
